\section{A Geometric Reference}
\label{sec:method}

The geometric reference tests whether wall crack locations drive association: it aligns wall surfaces to a common frame and compares projected crack masks.
We call this registration-gated Chamfer comparison \ours{} (query--gallery pairs whose images fail to register are not scored) and its registration-only variant, which assigns a constant score to registered pairs, \regonly{}.
The geometric reference assumes a roughly planar wall and uses no \textsc{CrackID} training.
Shared masks come from a segmentation model trained on public crack data only (no \textsc{CrackID} wall in training).
Crack instances use a minimum area of 200 pixels and a 5-pixel closing kernel, so all methods receive consistent inputs.
Masks are taken as given, and their error is not measured.

\subsection{Registration using wall texture}
\label{sec:registration}

For images $I_a$ and $I_b$, a homography $H$ ($x' \sim Hx$) is estimated from surrounding wall texture with cracks dilated and excluded, so alignment is driven by surface features independently of the tested identity.
This independence is conditional on the masks and on the cascade stage, not guaranteed per pair.
Full detector settings (SIFT with contrast-limited adaptive histogram equalization (CLAHE)~\cite{zuiderveld1994clahe}, ratio test, robust estimation) are archived with the code.

\subsection{Registration cascade and validity checks}
\label{sec:cascade}

Textured render generally provides background features, whereas smooth plaster does not.
The cascade has three stages: masked SIFT; unmasked SIFT if too few keypoints remain (flagged, since crack features can then influence alignment); and photometric ECC~\cite{evangelidis2008ecc}.
Candidate solutions are checked for plausibility: the quadrilateral must not be mirrored or folded, and the scale must lie within $[1/16,16]$; no minimum inlier count is imposed.
We record the successful stage per pair; its inlier root-mean-square reprojection error (RMSE) sets the comparison tolerance.
Absence of registration means lack of evidence, not that the cracks differ.

\subsection{Crack projection and geometric comparison}
\label{sec:chamfer}

Accepted registrations project instances from $I_a$ into $I_b$ and compare them using the symmetric mean Chamfer distance~\cite{barrow1977chamfer}, defined in~\eqref{eq:chamfer}. Let $a$ and $b$ be crack point sets and $d(p, b)=\min_{q \in b}\lVert p-q\rVert_2$:
\begin{equation}
D(a,b)=\frac{1}{2}\left[\frac{1}{|a|}\sum_{p \in a} d(p, b)+
\frac{1}{|b|}\sum_{q \in b} d(q, a)\right].
\label{eq:chamfer}
\end{equation}
Unlike intersection over union (IoU), which is overly sensitive to small misalignments of narrow cracks, Chamfer tolerates them. Candidates whose Chamfer distance exceeds the validation-calibrated threshold are rejected.
A scale-free alternative is Chamfer coverage, the fraction of crack points within the per-pair RMSE (reprojection error) of the other mask; tolerance is clamped to $[2,16]$~px, default $6$~px.

\subsection{Correspondence decisions and limits of spatial evidence}

For each image pair, one-to-one matching uses the Hungarian method~\cite{kuhn1955hungarian} at the validation-calibrated threshold; unregistered pairs score $-\infty$ (Section~\ref{sec:protocol}).
Alignment onto a crack-free region is evidence of absence; missed detections, out-of-frame areas, and registration failures leave the question unresolved.
The reference assumes that damage does not change; its performance may depend on the wall surface and on label provenance (Section~\ref{sec:annotation}).
